\documentclass[12pt]{article}
\usepackage[utf8]{inputenc}
\usepackage[T1]{fontenc}
\usepackage{hyperref}
\title{Themne Language Resources for K-12 Curriculum Development}
\author{Compiled by Claude Sonnet 5 (claude-sonnet-5) for Andr\'as Kornai}
\date{}
\begin{document}
\maketitle
\section{Introduction}
Themne (Temne; ISO 639-3 \texttt{tem}) is an Atlantic language of Sierra Leone. Search results showed Bible portions, a New Testament and literacy work, but no dictionary or school textbooks.
\section{Primary Educational Resources}
\begin{itemize}
\item Joshua Project lists a Themne literacy primer from Literacy \& Evangelism International: \url{https://www.joshuaproject.net/index.php/people_groups/15333/SL}
\item Lutheran Bible Translators, Themne Literacy and Scripture Engagement (reports over 30 literacy classes): \url{https://lbt.org/wp-content/uploads/2020/09/Themne-Literacy-and-Scripture-Engagement.pdf}
\end{itemize}
\section{Linguistic and Reference Materials}
\begin{itemize}
\item Wikipedia, Temne language (lists Kamarah, \emph{A descriptive grammar of Temne}, 2007, and Wilson, \emph{An outline of the Temne language}, 1961): \url{https://en.wikipedia.org/wiki/Temne_language}
\end{itemize}
\section{Religious and Community Resources}
\begin{itemize}
\item Lutheran Bible Translators, Sierra Leone (Old Testament drafted, checking pending): \url{https://lbt.org/country/sierra-leone}
\item Joshua Project, Temne-Banta in Sierra Leone: \url{https://joshuaproject.net/people_groups/10660/SL}
\end{itemize}
\section{Digital and Audio Resources}
None found in these searches.
\section{Curriculum Development Context}
The language is reported as having a New Testament (1868--1992 editions) but low literacy limits its use. A curriculum would build on the LBT literacy classes and the existing grammars; no primary-school textbooks were found.
\end{document}